\section{Proof Debugging: Turning Fog Into Work}
\chapterline

\subsection{The hidden scratchpad: why polished proofs hide the work}

In the preceding chapters, we previewed abstract algebra, topology, differential geometry, and algebraic geometry. But whatever the topic, doing mathematics eventually comes down to sitting with pencil and paper to write a proof. And often, progress stalls.

In mathematical textbooks and research papers, proofs are usually presented in final, polished form. The author omits the exploratory scratch work, dead ends, and tentative revisions that led to the result.

This presentation can make mathematical proof seem like a single uninterrupted deduction. When a beginner writes several lines, encounters an obstacle, and is unsure how to proceed, they may conclude that proof writing requires immediate insight rather than methodical troubleshooting.

In 1945, the Hungarian-American mathematician George P\'olya published \emph{How to Solve It}, emphasizing that mathematical problem-solving relies on systematic \emph{heuristics}---practical strategies and rules of thumb that guide work when no algorithmic recipe exists. Effective mathematicians rely on concrete methods: decomposing problems, testing minimal examples, and diagnosing why initial attempts fail.

\subsection{Treating proofs as structured arguments}

When a computer program crashes or produces the wrong output, a programmer does not stare at the screen waiting for inspiration; they open a debugger, inspect the variables, and find the exact line where the state went wrong.

Getting stuck on a proof works the same way. Hitting a wall usually points to a concrete issue---an unpacked definition, a missing hypothesis, a quantifier in the wrong order, or a false claim. Once you inspect a stalled draft the way you would debug code, getting stuck gives you something specific to check.

\begin{center}
\begin{tikzpicture}[
  node distance=1.1cm and 1.2cm,
  >=Latex,
  box/.style={draw=#1!80!black, fill=#1!12, thick, rounded corners=3pt, font=\small, align=center, inner sep=6pt},
  decision/.style={draw=quotientorange!80!black, fill=quotientorange!15, thick, rounded corners=3pt, font=\small\bfseries, align=center, inner sep=6pt},
  arrow/.style={->, thick, >=Latex, draw=gray!80}
]

  \node[box=trapred] (start) {\textbf{Stuck in the Fog?}\\Cannot see the next logical step};

  \node[decision, below=0.85cm of start] (example) {Test a Minimal Example\\($n=1, 2$, empty set, $2 \times 2$ matrix)};

  \node[box=trapred, right=1.4cm of example] (counter) {\textbf{Counterexample Found!}\\Claim is false or needs\\additional hypotheses};

  \node[box=sourceblue, below=0.85cm of example] (unpack) {\textbf{Unpack All Definitions}\\Translate every high-level term\\into its explicit quantifier formula};

  \node[box=structurepurple, below=0.85cm of unpack] (shelves) {\textbf{Separate Given vs. Goal}\\Given shelf: known facts and active variables\\Goal shelf: exact target shape to construct};

  \node[box=targetgreen, below=0.85cm of shelves] (match) {\textbf{Match Strategy to Goal}\\$\forall \implies$ choose arbitrary element\\$P \implies Q \implies$ assume $P$, deduce $Q$\\$\exists \implies$ construct explicit witness};

  \draw[arrow] (start) -- (example);
  \draw[arrow, trapred] (example) -- node[above, font=\scriptsize\bfseries, text=trapred] {Fails} (counter);
  \draw[arrow] (example) -- node[right, font=\scriptsize, text=gray!80!black] {Holds} (unpack);
  \draw[arrow] (unpack) -- (shelves);
  \draw[arrow] (shelves) -- (match);

\end{tikzpicture}\\[0.4em]
{\small\textbf{The Debugging Triage Flow:} When a proof draft stalls, test concrete instances before writing general logic. A counterexample immediately saves hours of futile effort; if the example holds, unpacking definitions into quantifier clauses reveals the structural moves required.}
\end{center}

\subsection{The five-stage diagnostic protocol}

When your pencil stops moving, execute these five steps in order:

\subsubsection{1. Test a minimal non-trivial example}
Before trying to prove an abstract theorem about arbitrary vector spaces, rings, or sequences, write down the smallest concrete instance you can think of. If a claim involves integers, test $n = 0, 1, 2, 3$. If it involves sets, test sets with $0$, $1$, and $2$ elements.
\begin{itemize}
\item If your claim fails on $n = 2$, you have found a counterexample, saving hours attempting to prove a false statement.
\item If your claim holds on small examples, observe \emph{why} it holds. The mechanism that makes $n = 2$ work often reveals the algebra needed for general $n$.
\end{itemize}

\textbf{Pause and think:} Suppose someone asserts: ``Every integer $n > 2$ can be written as the sum of two distinct prime numbers.'' Before trying an abstract proof, test minimal instances: $n = 3, 4, 5, 6$. For $n = 3$, the only positive integers smaller than $3$ are $1$ and $2$; since $1$ is not prime, $3$ cannot be written as a sum of two primes at all. For $n = 4$, the only prime smaller than $4$ is $2$, and $2 + 2 = 4$ repeats the same prime, so they are not distinct. The claim collapses on the first two instances you check! Testing small values prevents attempting to prove a falsehood.

\subsubsection{2. Unpack every definition down to its quantifier core}
Vague terminology obscures the underlying logic. A \emph{quantifier} is a logical prefix---such as $\forall$ (``for all'') or $\exists$ (``there exists'')---that explicitly declares how many elements satisfy a given condition. If the theorem says ``$f$ is continuous,'' do not leave the word ``continuous'' on your paper. Replace it with:
\[
\forall x_0 \in \R, \; \forall \eps > 0, \; \exists \delta > 0, \; \forall x \in \R, \; (|x - x_0| < \delta \implies |f(x) - f(x_0)| < \eps).
\]
In geometric terms, $\eps > 0$ specifies an error tolerance around the output $f(x_0)$, and $\delta > 0$ specifies a neighborhood radius around $x_0$ that guarantees function values remain within that tolerance. Once a definition is written in explicit logical symbols, the quantifiers determine the structure of the proof.

\subsubsection{3. Separate the Given shelf from the Goal shelf}
Take a blank sheet of paper and draw a horizontal line across the middle:
\begin{itemize}
\item \textbf{Top half (The Given Shelf)}: List every assumption, hypothesis, and active variable definition. These are the tools already in your hands.
\item \textbf{Bottom half (The Goal Shelf)}: State the exact target expression you must reach.
\end{itemize}
A common error is confusing an assumption with a target. Keeping them on separate physical halves of the paper prevents circular reasoning.

\begin{center}
\begin{tikzpicture}[scale=0.95, >=Latex]
  \draw[thick, sourceblue!80!black, fill=sourceblue!10, rounded corners=4pt]
    (-6.2, 1.8) rectangle (6.2, 3.8);
  \node[font=\bfseries, text=sourceblue!90!black, anchor=west] at (-5.8, 3.45) {The Given Shelf \normalfont(Tools in Your Hands)};
  \node[font=\small, align=left, anchor=west] at (-5.8, 2.65) {
    $\bullet$ Explicit hypotheses from the theorem statement\\
    $\bullet$ Active arbitrary variables introduced into scope ($\forall x \in S \implies x$ is fixed)\\
    $\bullet$ Unpacked definitions and proven intermediate facts
  };

  \draw[dashed, thick, quotientorange!80!black] (-4.0, 1.0) -- (-2.7, 1.0);
  \draw[dashed, thick, quotientorange!80!black] (2.7, 1.0) -- (4.0, 1.0);
  \node[fill=white, draw=quotientorange!80!black, thick, rounded corners=3pt, font=\small\bfseries, text=quotientorange!90!black, inner sep=4pt] at (0, 1.0) {Reasoning Bridge: Meet in the Middle};
  
  \draw[->, very thick, sourceblue!80!black] (-4.6, 1.65) -- (-4.6, 0.35) 
    node[midway, left=3pt, font=\scriptsize\bfseries, text=sourceblue!90!black, align=right] {Forward:\\derive consequences};
  \draw[->, very thick, targetgreen!80!black] (4.6, 0.35) -- (4.6, 1.65) 
    node[midway, right=3pt, font=\scriptsize\bfseries, text=targetgreen!90!black, align=left] {Backward:\\peel goal layers};

  \draw[thick, targetgreen!80!black, fill=targetgreen!10, rounded corners=4pt]
    (-6.2, -1.8) rectangle (6.2, 0.2);
  \node[font=\bfseries, text=targetgreen!90!black, anchor=west] at (-5.8, -0.15) {The Goal Shelf \normalfont(Target to Construct)};
  \node[font=\small, align=left, anchor=west] at (-5.8, -0.95) {
    $\bullet$ Outermost logical connective dictates your first proof sentence\\
    $\bullet$ Exact target shape to construct (e.g., find witness $w$, deduce $Q$)\\
    $\bullet$ Work backward: ``To show $Q$, it suffices to establish $P$''
  };
\end{tikzpicture}\\[0.4em]
{\small\textbf{The Two-Shelf Scratchpad:} By physically partitioning your scratch paper into a Given shelf (what is known) and a Goal shelf (what must be shown), you prevent circular reasoning. Work inward from both edges: peel quantifiers from the goal backward while extracting consequences from the givens forward until the two halves meet.}
\end{center}

\subsubsection{4. Type-check every mathematical symbol}
Just as in programming languages, every mathematical entity has a strict type:
\begin{itemize}
\item Is $x$ an element, a set, or a function?
\item Does the expression $f(A)$ mean the evaluation of a point or the image of a subset?
\item If $A$ and $B$ are sets, writing $A < B$ is a type error unless an order relation has been explicitly defined on the power set.
\end{itemize}
Type mismatches account for more than half of all beginner proof errors.

\textbf{Pause and think:} Suppose $f: X \to Y$ is a function between sets, and $B \subseteq Y$. What is the type of $f^{-1}(B)$ versus $f^{-1}(y)$ for a point $y \in Y$? The preimage $f^{-1}(B) = \{x \in X : f(x) \in B\}$ is always a well-defined subset of $X$, regardless of whether $f$ has an inverse. But writing $f^{-1}(y)$ treats $f^{-1}$ as an inverse function mapping a point to a point, which is a type error unless $f$ is already known to be bijective (a one-to-one matching). Confusing a set with an element is one of the most common compilation errors in beginner proofs.

\subsubsection{5. Verify the connective architecture}
Look at the outermost logical connective of your goal:
\begin{itemize}
\item If the goal is $\forall x \in S, \, P(x)$, your proof \emph{must} begin: ``Let $x \in S$ be arbitrary.''
\item If the goal is $P \implies Q$, your proof \emph{must} begin: ``Assume $P$.'' Your temporary goal becomes $Q$.
\item If the goal is $A \subseteq B$, your proof \emph{must} begin: ``Let $x \in A$.'' Your target is to show $x \in B$.
\item If the goal is $\exists x \in S, \, P(x)$, your proof must name or construct a specific witness $x \in S$ and demonstrate that $P(x)$ holds.
\item If the goal is $P \land Q$, your proof must split into two independent sub-proofs: first deduce $P$, then deduce $Q$.
\item If the goal is $P \lor Q$, you may assume $\neg P$ and deduce $Q$ (since $P \lor Q$ is logically equivalent to $\neg P \implies Q$), or divide your argument into exhaustive cases.
\end{itemize}
The grammar of the goal dictates the first line of your proof before you even understand the tricky mathematics inside.

\begin{center}
\begin{tikzpicture}[scale=0.92, >=Latex]
  % Row 1: Universal
  \node[draw=structurepurple!80!black, fill=structurepurple!10, rounded corners=3pt, thick, text width=4.8cm, align=center, font=\small] (g1) at (-4.2, 2.4) {
    \textbf{Goal Form: Universal}\\
    $\forall x \in S, \, P(x)$
  };
  \node[draw=targetgreen!80!black, fill=targetgreen!10, rounded corners=3pt, thick, text width=6.6cm, align=left, font=\small] (s1) at (3.2, 2.4) {
    \textbf{Mandated Opening Move:}\\
    ``Let $x \in S$ be arbitrary.'' \hfill (\emph{target: $P(x)$})
  };
  \draw[->, very thick, gray!70] (g1) -- (s1);

  % Row 2: Implication
  \node[draw=structurepurple!80!black, fill=structurepurple!10, rounded corners=3pt, thick, text width=4.8cm, align=center, font=\small] (g2) at (-4.2, 0.8) {
    \textbf{Goal Form: Implication}\\
    $P \implies Q$
  };
  \node[draw=targetgreen!80!black, fill=targetgreen!10, rounded corners=3pt, thick, text width=6.6cm, align=left, font=\small] (s2) at (3.2, 0.8) {
    \textbf{Mandated Opening Move:}\\
    ``Assume $P$.'' \hfill (\emph{new target: $Q$})
  };
  \draw[->, very thick, gray!70] (g2) -- (s2);

  % Row 3: Subset
  \node[draw=structurepurple!80!black, fill=structurepurple!10, rounded corners=3pt, thick, text width=4.8cm, align=center, font=\small] (g3) at (-4.2, -0.8) {
    \textbf{Goal Form: Set Inclusion}\\
    $A \subseteq B$
  };
  \node[draw=targetgreen!80!black, fill=targetgreen!10, rounded corners=3pt, thick, text width=6.6cm, align=left, font=\small] (s3) at (3.2, -0.8) {
    \textbf{Mandated Opening Move:}\\
    ``Let $x \in A$.'' \hfill (\emph{new target: $x \in B$})
  };
  \draw[->, very thick, gray!70] (g3) -- (s3);

  % Row 4: Existential
  \node[draw=structurepurple!80!black, fill=structurepurple!10, rounded corners=3pt, thick, text width=4.8cm, align=center, font=\small] (g4) at (-4.2, -2.4) {
    \textbf{Goal Form: Existential}\\
    $\exists x \in S, \, P(x)$
  };
  \node[draw=targetgreen!80!black, fill=targetgreen!10, rounded corners=3pt, thick, text width=6.6cm, align=left, font=\small] (s4) at (3.2, -2.4) {
    \textbf{Mandated Opening Move:}\\
    ``Construct/name candidate witness $x \in S$.''
  };
  \draw[->, very thick, gray!70] (g4) -- (s4);
\end{tikzpicture}\\[0.4em]
{\small\textbf{Connective Architecture:} The outermost logical connective of the goal strictly dictates the opening line of your proof.}
\end{center}

\subsection{Catalog of proof gaps and repairs}

\begin{longtable}{p{0.18\linewidth}p{0.35\linewidth}p{0.39\linewidth}}
\toprule
\textbf{Gap Name} & \textbf{Symptom in Draft} & \textbf{Surgical Repair} \\
\midrule
\textbf{Quantifier Gap} & An object is manipulated without introducing where it came from or what role it plays. & Declare the variable explicitly: state whether it is arbitrary ($\forall$) or chosen ($\exists$), and specify its domain. \\
\addlinespace
\textbf{Definition Gap} & Words like ``compact'' or ``divisible'' are treated intuitively without invoking their mathematical equations. & Erase the English word and replace it with its exact definition. \\
\addlinespace
\textbf{Existence Gap} & A proof asserts ``choose an element $x$ with property $P$'' when no such element is guaranteed to exist. & Construct an explicit formula for $x$, or cite an established theorem (like the Intermediate Value Theorem) that guarantees existence. \\
\addlinespace
\textbf{Implication Gap} & A line jumps directly from step $A$ to step $B$ with the word ``clearly,'' hiding a non-trivial theorem. & Break the leap into two smaller steps; write down the intermediate lemma and verify it. \\
\addlinespace
\textbf{Case Gap} & An algebraic step divides by $x - 1$ without considering the possibility that $x = 1$. & Partition the universe into explicit exhaustive cases: Case 1 ($x = 1$) and Case 2 ($x \neq 1$). \\
\addlinespace
\textbf{Converse Error} & Attempting to prove $P \implies Q$ by assuming $Q$ and showing it leads to $P$. & Reorient the proof: start strictly with hypothesis $P$ and verify that each arrow points forward toward $Q$. \\
\bottomrule
\end{longtable}

\subsection{Case study: debugging a broken parity proof}

Consider the following student proof attempt regarding \emph{parity} (whether an integer is even or odd):

\begin{quote}
\textbf{Claim}: If $n^2$ is even, then $n$ is even. \\
\textbf{Draft Proof}: Since $n^2$ is even, $n^2 = 2k$ for some integer $k$. Taking square roots of both sides gives $n = \sqrt{2k}$. Therefore $n$ is even. \hfill $\square$
\end{quote}

\textbf{Debugging Analysis}:
\begin{enumerate}
\item \emph{Line 1}: $n^2 = 2k$ for $k \in \Z$. Valid; correctly unpacks the definition of even.
\item \emph{Line 2}: $n = \sqrt{2k}$. Valid algebraic step (assuming $n \ge 0$).
\item \emph{Line 3}: ``Therefore $n$ is even.'' \textbf{Bug detected!} The definition of even requires $n = 2m$ for an integer $m$. Writing $n = \sqrt{2k} = \sqrt{2}\sqrt{k}$ places the factor of $2$ inside a radical, which does not directly establish that $n$ is an integer multiple of $2$. To deduce from $\sqrt{2k}$ that $n$ is an even integer requires showing that $k = 2m^2$, which assumes the conclusion. Taking square roots moves the argument out of integer arithmetic into the real numbers.
\end{enumerate}

\textbf{Surgical Repair via Contrapositive}:
Taking square roots is inconvenient because square roots do not preserve integer divisibility directly. Instead, consider the \emph{contrapositive} (proving ``if $P$ then $Q$'' by establishing the logically equivalent statement ``if not $Q$ then not $P$''):
\[
\text{If } n \text{ is odd, then } n^2 \text{ is odd.}
\]
Now the algebra works forward smoothly:
\begin{proof}
We prove the contrapositive. Assume $n$ is odd. By definition, there exists an integer $m$ such that $n = 2m + 1$. Squaring both sides yields:
\[
n^2 = (2m + 1)^2 = 4m^2 + 4m + 1 = 2(2m^2 + 2m) + 1.
\]
Since $m$ is an integer, $k = 2m^2 + 2m$ is an integer. Thus $n^2 = 2k + 1$, which means $n^2$ is odd. Having proved the contrapositive, the original claim holds.
\end{proof}

\begin{center}
\begin{tikzpicture}[scale=0.92, >=Latex]
  % Direct Track (Top)
  \node[font=\footnotesize\bfseries, trapred!90!black, anchor=west] at (-6.3, 1.8) {Direct attempt ($P \implies Q$):};
  
  \node[draw=trapred!70!black, fill=trapred!8, rounded corners=3pt, thick, text width=3.1cm, align=center, font=\scriptsize] (d1) at (-4.5, 0.8) {
    \textbf{1. Unpack given}\\
    $n^2 = 2k$ for $k \in \Z$
  };
  \node[draw=trapred!70!black, fill=trapred!8, rounded corners=3pt, thick, text width=3.1cm, align=center, font=\scriptsize] (d2) at (-0.6, 0.8) {
    \textbf{2. Solve for $n$}\\
    $n = \sqrt{2k} = \sqrt{2}\sqrt{k}$
  };
  \node[draw=trapred!90!black, fill=trapred!20, rounded corners=3pt, thick, text width=4.1cm, align=center, font=\scriptsize, inner sep=4pt] (d3) at (3.9, 0.8) {
    \textbf{\textcolor{trapred!90!black}{DEAD END}}\\[0.3em]
    Factor $2$ trapped: $\sqrt{2}$ irrational;\\[0.3em]
    cannot extract integer $2$\\[0.2em]
    without circularity!
  };
  \draw[->, thick, trapred] (d1) -- (d2);
  \draw[->, thick, trapred] (d2) -- (d3);

  % Equivalence indicator
  \draw[<->, very thick, quotientorange] (-6.0, 0.3) -- (-6.0, -0.9)
    node[midway, right=2pt, font=\tiny\bfseries, text=quotientorange!90!black] {Contrapositive switch};

  % Contrapositive Track (Bottom)
  \node[font=\footnotesize\bfseries, targetgreen!90!black, anchor=west] at (-6.3, -1.1) {Contrapositive path ($\neg Q \implies \neg P$):};
  
  \node[draw=targetgreen!80!black, fill=targetgreen!8, rounded corners=3pt, thick, text width=3.1cm, align=center, font=\scriptsize, inner sep=4pt] (c1) at (-4.5, -2.1) {
    \textbf{1. Assume $\neg Q$}\\
    $n$ is odd: $n = 2m+1$
  };
  \node[draw=targetgreen!80!black, fill=targetgreen!8, rounded corners=3pt, thick, text width=3.1cm, align=center, font=\scriptsize, inner sep=4pt] (c2) at (-0.6, -2.1) {
    \textbf{2. Compute $n^2$}\\
    $n^2 = 4m^2 + 4m + 1$
  };
  \node[draw=targetgreen!90!black, fill=targetgreen!20, rounded corners=3pt, thick, text width=4.1cm, align=center, font=\scriptsize, inner sep=4pt] (c3) at (3.9, -2.1) {
    \textbf{\textcolor{targetgreen!90!black}{SUCCESS}}\\[0.3em]
    Integer witness: $j = 2m^2+2m \in \Z$;\\[0.3em]
    $n^2 = 2j+1 \implies n^2$ is odd!
  };
  \draw[->, thick, targetgreen] (c1) -- (c2);
  \draw[->, thick, targetgreen] (c2) -- (c3);
\end{tikzpicture}\\[0.4em]
{\small\textbf{Proof Debugging in Action:} Direct algebraic manipulation of $n^2=2k$ traps the factor of $2$ inside an irrational radical $\sqrt{2k}$. Switching to the contrapositive ($n \text{ odd} \implies n^2 \text{ odd}$) works exclusively within polynomial integer arithmetic, revealing an explicit integer witness.}
\end{center}

\begin{labbox}
Keep a dedicated ``bug log'' in your notebook. Whenever your proof stalls or an instructor finds an unsupported step, write down the name of the gap (from the table above) and the exact correction you applied. Patterns will emerge rapidly.
\end{labbox}

\begin{historybox}
In his 1976 study \emph{Proofs and Refutations}, Imre Lakatos argued that mathematical knowledge often develops through an iterative process of conjectures, counterexamples, and conceptual refinements.

Lakatos reconstructed the historical development of Euler's formula for polyhedra---three-dimensional solid shapes bounded by flat polygonal faces, straight edges, and sharp vertices. In 1758, Leonhard Euler observed that for familiar solids like the cube and tetrahedron, the count of vertices $V$, edges $E$, and faces $F$ obeys $V - E + F = 2$. In 1813, Augustin-Louis Cauchy published an argument: remove one face, stretch the remaining network flat onto a plane, and triangulate each polygon. Yet over the following century, counterexamples emerged. Self-intersecting star polyhedra with interlocking vertices violated the count. A small cube nested hollowly inside a larger cube gave $V - E + F = 4$, and a picture-frame polyhedron with a central tunnel gave $V - E + F = 0$.

When confronted with these counterexamples, mathematicians initially debated whether to exclude non-standard polyhedra (an approach Lakatos called \emph{monster-barring}) or to modify definitions to accommodate them. The mathematical resolution came from inspecting where Cauchy's argument failed: Cauchy had assumed that every polyhedron could be inflated into a simple sphere without holes. By identifying that topological assumption and incorporating it as an explicit hypothesis, mathematicians proved the general relation: $V - E + F = 2 - 2g$, where $g$ is the topological genus (the number of doughnut-like tunnels passing through the surface; $g = 0$ for a sphere, while $g = 1$ for a picture frame).

In modern mathematics and computer science, proof checking has extended to formal interactive theorem provers such as Lean, Coq, and Isabelle. While these systems verify whether a formalised proof script is completely verified, an active area of research is automated proof repair: given an informal, incomplete, or buggy proof draft, how can an algorithm isolate the missing lemma, find a counterexample, or distinguish between a notation error and a conceptual gap?
\end{historybox}

\subsection{The craft of proof debugging---and the practice ahead}

Getting stuck in the middle of a proof is normal; what matters is having concrete moves to make next. The five-stage diagnostic protocol---testing minimal examples, unpacking definitions into explicit quantifiers, separating the Given and Goal shelves, type-checking every symbol, and matching the proof structure to the main connective---turns a vague sense of confusion into specific checks. Even when a draft fails outright, the counterexample that breaks it shows which hypothesis was missing or why the claim needs to be refined.

Diagnostic protocols become reliable habits through regular practice on concrete problems. To build confidence and fluency before reading standard undergraduate texts, Chapter~26 presents \textbf{A Serious Velleman Readiness Plan}, providing a four-week practice progression and a ten-proof benchmark portfolio.

\subsection{Exercises}

\begin{enumerate}
\item \textbf{Mental Check (Spot the Gap Type):} For each of the following flawed proof fragments, identify the primary gap from the catalog (Quantifier Gap, Definition Gap, Existence Gap, Implication Gap, Case Gap, or Converse Error):
\begin{enumerate}[label=(\alph*)]
\item ``To prove that $f$ is continuous on $\R$, suppose $|x - y|$ is very small. Then $f(x)$ is close to $f(y)$, so $f$ is continuous.''
\item ``We wish to prove that if $n$ is a multiple of $6$, then $n$ is even. Assume $n$ is even, so $n = 2k$. Choosing $k = 3m$ gives $n = 6m$, so $n$ is a multiple of $6$.''
\item ``Let $x^2 - x = 0$. Dividing both sides by $x$ gives $x - 1 = 0$, so $x = 1$. Therefore $1$ is the unique solution.''
\end{enumerate}

\item \textbf{Routine Practice (Debugging Parity and Divisibility):} A student attempts to prove the following claim:
\begin{quote}
\textbf{Claim:} For all integers $n$, if $3n + 2$ is odd, then $n$ is odd.\\
\textbf{Student Draft:} Assume $3n + 2$ is odd. Then $3n + 2 = 2k + 1$ for some integer $k$. Subtracting $2$ yields $3n = 2k - 1$. Dividing by $3$ gives $n = \frac{2k - 1}{3}$. Since $2k - 1$ is odd, $n$ must be odd. \hfill $\square$
\end{quote}
\begin{enumerate}[label=(\alph*)]
\item Diagnose the fatal bug in the student's draft. (Why does $\frac{2k - 1}{3}$ fail to prove that $n$ is an odd integer?)
\item Formulate the contrapositive of the claim.
\item Write a complete, correct proof using the contrapositive.
\end{enumerate}

\item \textbf{Routine Practice (Element Chasing and Case Splitting):} A student writes the following proof that $A \cap (B \cup C) \subseteq (A \cap B) \cup C$ for arbitrary sets $A, B, C$:
\begin{quote}
\textbf{Student Draft:} Let $x \in A \cap (B \cup C)$. Then $x \in A$ and $x \in B \cup C$. Thus $x \in A$ and $x \in B$. Therefore $x \in A \cap B$, which implies $x \in (A \cap B) \cup C$. \hfill $\square$
\end{quote}
\begin{enumerate}[label=(\alph*)]
\item Point out the exact line where the draft makes an unjustified leap.
\item Provide a concrete counterexample showing that $x \in B \cup C$ does not imply $x \in B$.
\item Supply the missing repair by partitioning into cases ($x \in B$ versus $x \in C$) to complete a valid proof.
\end{enumerate}

\item \textbf{Deeper Inquiry (The All-Horses Fallacy):} Consider the famous false theorem: ``All horses have the same color.'' The induction attempt argues:
\begin{itemize}
\item \emph{Base step ($n = 1$):} In any set consisting of one horse, all horses in the set trivially have the same color.
\item \emph{Induction step:} Assume any set of $k$ horses has the same color. Given a set of $k+1$ horses $\{h_1, h_2, \dots, h_{k+1}\}$, the first $k$ horses $\{h_1, \dots, h_k\}$ share the same color by the induction hypothesis. Similarly, the last $k$ horses $\{h_2, \dots, h_{k+1}\}$ share the same color. Since the two groups overlap at the middle horses, all $k+1$ horses must share the same color.
\end{itemize}
Test this argument on minimal examples ($k = 1, 2$). For which value of $k$ does the deduction bridge collapse, and what exact set-theoretic condition fails?

\item \textbf{Deeper Inquiry (Preserving Intersections Under Maps):} Consider the claim that for any function $f: X \to Y$ and any subsets $A, B \subseteq X$:
\[
f(A \cap B) = f(A) \cap f(B).
\]
\begin{enumerate}[label=(\alph*)]
\item Test a minimal finite example: let $X = \{1, 2\}$, $Y = \{0\}$, and let $f(1) = f(2) = 0$. Choose $A = \{1\}$ and $B = \{2\}$. Compute $f(A \cap B)$ and $f(A) \cap f(B)$. Is the claim true?
\item Set up the Two-Shelf Scratchpad to prove the containment $f(A \cap B) \subseteq f(A) \cap f(B)$. Write a complete, rigorous proof of this inclusion.
\item What property must the function $f$ possess (e.g., injectivity or surjectivity) to guarantee the reverse inclusion $f(A) \cap f(B) \subseteq f(A \cap B)$? Prove your assertion.
\end{enumerate}
\end{enumerate}

\newpage
